\documentclass[11pt,a4paper]{article}   % 11 pt text on A4 paper

% Page
\usepackage[margin=2.5cm]{geometry}     % 2.5 cm margins all round
\usepackage{parskip}                    % a blank line between paragraphs instead of an indent
\usepackage[T1]{fontenc}                % real < > | " characters in code, and text that copies cleanly from the PDF
\usepackage{lmodern}                    % Latin Modern, a crisp scalable version of the default font

% Maths
\usepackage{amsmath}                    % align, cases and \text inside maths

% Figures and tables
\usepackage{booktabs}                   % clean table rules: \toprule, \midrule, \bottomrule

% Code
\usepackage{xcolor}                     % named colours for the code highlighting below
\usepackage{listings}                   % highlighted code blocks and \lstinline, with no shell-escape
\usepackage{algorithm}                  % numbered, captioned boxes for algorithms
\usepackage{algpseudocode}              % pseudocode commands: \State, \If, \While, \Function, ...
\algrenewcommand\algorithmicrequire{\textbf{Input:}}    % \Require prints Input: rather than Require:
\algrenewcommand\algorithmicensure{\textbf{Output:}}    % \Ensure prints Output: rather than Ensure:

% How code looks: set language=... per listing (Python, C, Java, Matlab, bash, SQL, ...).
% Under pdfLaTeX listings only reads plain ASCII, so keep accents and emoji out of code.
\lstdefinestyle{report}{
  basicstyle       = \ttfamily\small,
  keywordstyle     = \color{blue!70!black}\bfseries,
  commentstyle     = \color{green!40!black}\itshape,
  stringstyle      = \color{red!60!black},
  numbers          = left,              % line numbers down the left
  numberstyle      = \tiny\color{gray},
  frame            = single,
  rulecolor        = \color{gray!40},
  backgroundcolor  = \color{gray!5},
  breaklines       = true,              % wrap long lines instead of running off the page
  postbreak        = \mbox{\textcolor{gray}{$\hookrightarrow$}\space},
  columns          = fullflexible,      % copying from the PDF gives back the code as written
  keepspaces       = true,
  showstringspaces = false,
  upquote          = true,              % straight quotes, as typed
  tabsize          = 4,
  xleftmargin      = 1.5em,             % room for the line numbers
  framexleftmargin = 0.5em
}
\lstset{style=report}

% Links
\usepackage[hidelinks]{hyperref}        % clickable cross-references, without coloured boxes
\hypersetup{                            % details shown in the PDF viewer’s document properties
  pdftitle  = {Assignment 2: Shortest paths in weighted graphs},
  pdfauthor = {Your Name},
  pdflang   = {en-AU}
}

\title{Assignment 2: Shortest paths in weighted graphs}
\author{Your Name\\ Student number 12345678}
\date{Due 24 October 2026}

\begin{document}
\maketitle

\section{Introduction}

This report describes an implementation of Dijkstra's algorithm, which finds the length of the
shortest path from one vertex to every other vertex of a graph whose edge weights are
non-negative. Section~\ref{sec:method} gives the algorithm, Section~\ref{sec:code} the code, and
Section~\ref{sec:analysis} its running time.

\section{Method}\label{sec:method}

Algorithm~\ref{alg:dijkstra} keeps a tentative distance $d[v]$ for every vertex and a
priority queue of vertices to visit, always taking the closest unvisited vertex next. Rather than
lowering a key already in the queue, it pushes a new entry and skips stale ones when they come out,
which suits a plain binary heap.

\begin{algorithm}[ht]
  \caption{Dijkstra's algorithm with a binary heap}
  \label{alg:dijkstra}
  \begin{algorithmic}[1]
    \Require a graph $G = (V, E)$ with weights $w(u, v) \ge 0$, and a source $s \in V$
    \Ensure $d[v]$, the length of a shortest path from $s$ to each $v$
    \State $d[v] \gets \infty$ for every $v \in V$, and $d[s] \gets 0$
    \State $Q \gets \{(0, s)\}$ \Comment{a min-priority queue of (distance, vertex)}
    \While{$Q$ is not empty}
      \State $(k, u) \gets \Call{ExtractMin}{Q}$
      \If{$k > d[u]$}
        \State \textbf{continue} \Comment{a stale entry}
      \EndIf
      \ForAll{edges $(u, v) \in E$}
        \If{$d[u] + w(u, v) < d[v]$}
          \State $d[v] \gets d[u] + w(u, v)$
          \State \Call{Insert}{$Q, (d[v], v)$}
        \EndIf
      \EndFor
    \EndWhile
    \State \Return $d$
  \end{algorithmic}
\end{algorithm}

\section{Implementation}\label{sec:code}

Listing~\ref{lst:python} implements Algorithm~\ref{alg:dijkstra} in Python, using the standard
library's \lstinline[language=Python]|heapq| module as the queue. The call
\lstinline[language=Python]|heapq.heappop(queue)| returns the entry with the smallest distance.

% float=tp moves the listing to the top of a page rather than splitting it; leave it off for code longer than a page
\begin{lstlisting}[language=Python, float=tp, caption={Dijkstra's algorithm in Python.}, label={lst:python}]
import heapq

def dijkstra(graph, source):
    """Return the shortest distance from source to every reachable vertex."""
    dist = {source: 0}
    queue = [(0, source)]
    while queue:
        d, u = heapq.heappop(queue)
        if d > dist[u]:
            continue  # stale entry: u was already reached by a shorter path
        for v, weight in graph[u]:
            if d + weight < dist.get(v, float("inf")):
                dist[v] = d + weight
                heapq.heappush(queue, (dist[v], v))
    return dist

graph = {
    "A": [("B", 4), ("C", 1)],
    "B": [("D", 1)],
    "C": [("B", 2), ("D", 5)],
    "D": [],
}
print(dijkstra(graph, "A"))  # {'A': 0, 'B': 3, 'C': 1, 'D': 4}
\end{lstlisting}

The inner step, relaxing one edge, is the same in any language. Listing~\ref{lst:c} shows it in C,
where an unreached vertex has distance \lstinline[language=C]|INT_MAX|.

\begin{lstlisting}[language=C, caption={Relaxing an edge in C.}, label={lst:c}]
#include <limits.h>

/* Relax the edge (u, v) with weight w; returns 1 when dist[v] got shorter. */
int relax(int dist[], int u, int v, int w)
{
    if (dist[u] != INT_MAX && dist[u] + w < dist[v]) {
        dist[v] = dist[u] + w;
        return 1;
    }
    return 0;
}
\end{lstlisting}

\section{Analysis}\label{sec:analysis}

Every edge is relaxed at most once, and each relaxation can push one entry onto the queue, so
with a binary heap the algorithm runs in $O\bigl((|V| + |E|) \log |V|\bigr)$ time.
Table~\ref{tab:queues} compares this with other choices of priority queue.

\begin{table}[ht]
  \centering
  \caption{Running time of Dijkstra's algorithm by priority queue.}
  \label{tab:queues}
  \begin{tabular}{lccc}
    \toprule
    Priority queue & Decrease-key & Extract-min & Total \\
    \midrule
    Unsorted array & $O(1)$ & $O(|V|)$ & $O(|V|^2)$ \\
    Binary heap & $O(\log |V|)$ & $O(\log |V|)$ & $O\bigl((|V| + |E|) \log |V|\bigr)$ \\
    Fibonacci heap & $O(1)$ amortised & $O(\log |V|)$ amortised & $O(|E| + |V| \log |V|)$ \\
    \bottomrule
  \end{tabular}
\end{table}

For sparse graphs, where $|E|$ is close to $|V|$, the binary heap is the practical choice; the
unsorted array only wins on dense graphs with $|E|$ close to $|V|^2$.

\end{document}
